\begin{env}[5.13.5]
\label{IV.2.s5.13.5-fr}
Diremos que un anillo $A$ es \emph{normal} si
$\operatorname{Spec}(A)$ es un esquema \emph{normal}, dicho de otro modo, si para
todo ideal primo $\mathfrak p$ de $A$, $A_{\mathfrak p}$ es \emph{íntegro e
integralmente cerrado}. Se observará que esto implica que $A$ es \emph{reducido} y que
$\mathfrak p$ solo puede contener \emph{un solo} ideal primo minimal $\mathfrak q$
de $A$, dicho de otro modo, que $\mathfrak p$ pertenece a \emph{una sola} componente
irreducible de $\operatorname{Spec}(A)$; además, $\mathfrak q$ es la imagen
inversa de $(0)$ por el homomorfismo canónico $A\to A_{\mathfrak p}$, y
$A_{\mathfrak p}$ es, por tanto, isomorfo a
$(A/\mathfrak q)_{\mathfrak p/\mathfrak q}$, lo que prueba que el anillo
íntegro $A/\mathfrak q$ es \emph{integralmente cerrado} como intersección de anillos
integralmente cerrados (\textbf I, 8.2.1.1). Si $A$ es \emph{noetheriano} y \emph{normal}, $A$ es
producto directo de un número finito de anillos íntegros e integralmente
cerrados (\textbf{I}, 5.1.4).
\end{env}

\begin{corollary}[5.13.6]
\label{IV.2.s5.13.6-fr}
Supongamos que los $A_\beta$ son normales y que, para
$\alpha\leq\beta$, toda componente irreducible de
$\operatorname{Spec}(A_\beta)$ domina una componente irreducible de
$\operatorname{Spec}(A_\alpha)$. Entonces $A=\varinjlim A_\alpha$ es normal.
\end{corollary}

En efecto, sea $\mathfrak p$ un ideal primo de $A$ y pongamos, para todo
$\alpha$, $\mathfrak p_\alpha=\varphi_\alpha^{-1}(\mathfrak p)$. Por
hipótesis, para todo $\alpha$, $\mathfrak p_\alpha$ contiene un único ideal
minimal $\mathfrak q_\alpha$ de $A_\alpha$, y para $\alpha\leq\beta$,
$\varphi_{\beta\alpha}^{-1}(\mathfrak q_\beta)$ es un ideal minimal de
$A_\alpha$ contenido en $\mathfrak p_\alpha$, por tanto igual a
$\mathfrak q_\alpha$. Si se pone
$B_\alpha=A_\alpha/\mathfrak q_\alpha$, los anillos íntegros $B_\alpha$
forman, pues, un sistema inductivo para los homomorfismos
$\varphi'_{\beta\alpha}:B_\alpha\to B_\beta$ deducidos de los
$\varphi_{\beta\alpha}$ por paso a los cocientes, y que son inyectivos por
definición; además, los anillos $B_\alpha$ son integralmente cerrados, por
tanto lo mismo ocurre con $B=\varinjlim B_\alpha$ (5.13.4). Si se pone
$\mathfrak p'_\alpha=\mathfrak p_\alpha/\mathfrak q_\alpha$, se tiene además
$(B_\alpha)_{\mathfrak p'_\alpha}=(A_\alpha)_{\mathfrak p_\alpha}$; como
$\varinjlim(B_\alpha)_{\mathfrak p'_\alpha}
=\varinjlim(A_\alpha)_{\mathfrak p_\alpha}=A_{\mathfrak p}$
es un anillo local de $B$ en un ideal primo (5.13.3), es un anillo íntegro e
integralmente cerrado, lo que demuestra el corolario.
